\documentclass[11pt]{article}
\usepackage[margin=1in]{geometry}
\usepackage{amsmath}
\usepackage{microtype}
\usepackage[hidelinks]{hyperref}

\newcommand{\Faq}{\textit{F.~aquilonia}}
\newcommand{\Fpo}{\textit{F.~polyctena}}
\newcommand{\BF}{\ensuremath{\mathrm{BF}(\mathrm{dB})}}

\title{\vspace{-3em}Ancestry-diagnostic loci are enriched in climate-associated
regions,\\ but directional sorting is not}
\author{}
\date{}

\begin{document}
\maketitle
\vspace{-4em}

\noindent\textit{Draft results summary (Module C, the extrinsic-climate arm).
Methods are deferred to Supplementary Materials. Findings are descriptive:
climate associations are genotype--environment correlations, and the sorting
itself is licensed as selection only against the forthcoming neutral null. The
identity of the climate axes (which variables load on PC1 versus PC2) is
pending.}

\section*{Results}

\paragraph{Design.}
Per-population allele frequencies were associated with two per-population climate
principal components (PC1, PC2) using BayPass across eight configurations: two
population sets, two climate axes, and two covariance settings. Because sorting is
by definition ancestry fixation, a climate axis that is itself correlated with
genome-wide ancestry could co-localise with sorting for structural rather than
ecological reasons. PC1 is only weakly ancestry-correlated across populations
(Spearman $\rho = -0.23$); PC2 is moderately so ($\rho = +0.57$), driven
predominantly by a single population, {\AA}land. We therefore condition on the
among-population covariance matrix estimated from the LD-pruned marker set, and
take the {\AA}land-excluded population set as the primary analysis, with the
remaining configurations as comparators.

\paragraph{Measuring climate association per linkage cluster.}
We summarise each LD cluster's climate association by its \emph{significance
rate}: the proportion of its member loci reaching $\BF \ge 15$. A rate is used
rather than a count of significant members because a count makes a cluster's
eligibility depend on its size---a cluster of nine loci (the median) would need
56\% of its members significant to reach a threshold of five, whereas a cluster of
ninety needs only 6\%---so a count-based criterion applies a far stricter standard
to small clusters and cannot be corrected by conditioning on size alone. All tests
treat the cluster as the unit of observation, since member loci within a cluster
are in linkage disequilibrium and are not independent; cluster size is retained as
a covariate throughout. Because a rate is estimated imprecisely when a cluster
contains few loci, results are reported across nested size strata.

\paragraph{Ancestry-diagnostic loci concentrate in climate-associated regions.}
Among clusters large enough for the significance rate to be well estimated
($\ge 50$ loci), the proportion of ancestry-diagnostic members (DI $> -25$)
increases with climate association: each $+10$ percentage points of significance
rate corresponds to $+4.5$ percentage points of diagnostic loci for PC1
($p = 8\times10^{-4}$) and $+3.7$ for PC2 ($p = 0.013$), against a baseline of
roughly 10\%. A size-weighted, overdispersion-corrected model gives concordant
estimates (odds ratios 1.50 and 1.47 per $+10$ percentage points). The effect is
attenuated in the intermediate stratum ($\ge 20$ loci: $+1.9$ and $+1.8$
percentage points) and is not detectable when all clusters are pooled, as expected
when the predictor is measured with substantial error in small clusters. Both
climate axes behave alike, at similar magnitude.

\paragraph{Directional sorting shows no such association.}
Applying the same size-normalised, cluster-level test to directional ancestry
sorting yields no relationship in the size-comparable stratum (odds ratios 0.99
for PC1 and 1.03 for PC2, $p > 0.95$); a member-level analysis is likewise null.
This is the expected result on mechanistic grounds: genotype--environment
association is detected through variation in allele frequency \emph{among}
populations, whereas unidirectional sorting drives most populations toward the
same near-fixed state and therefore minimises exactly that variance. Climate-driven
selection would in any case predict populations diverging in different directions
according to local conditions, rather than the parallel fixation that
unidirectional sorting describes.

\paragraph{Interpretation and outlook.}
The climate signal in these data is a \emph{static} one: regions whose allele
frequencies track climate are enriched for ancestry-diagnostic loci, on both
climate axes, by a modest but consistent margin. There is no corresponding
\emph{dynamic} signal---directional ancestry sorting is not preferentially located
in climate-associated regions. We note the limits: the informative stratum
comprises 878 clusters, the enrichment is a few percentage points against a
variable background, and these are associations rather than demonstrations of
selection. Identifying the climate axes and evaluating the sorting against a
recombination-matched neutral null remain the decisive next steps; intrinsic
incompatibilities are pursued separately through among-population linkage
disequilibrium between unlinked diagnostic units.

\end{document}
